% =====================================================================
%  Beamer Presentation — 10 slides following script.tex exactly
%  Experimental Analysis of Sorting Algorithms
%  Maria Maiorescu -- West University of Timișoara
%  Compile: pdflatex presentation.tex  (run twice)
% =====================================================================
\documentclass[aspectratio=169]{beamer}

% --- Theme ---
\usetheme{Madrid}
\usecolortheme{whale}
\setbeamertemplate{navigation symbols}{}
\setbeamertemplate{footline}{
  \leavevmode%
  \hbox{%
    \begin{beamercolorbox}[wd=.33\paperwidth,ht=2.25ex,dp=1ex,center]{author in head/foot}%
      \usebeamerfont{author in head/foot}Maria Maiorescu
    \end{beamercolorbox}%
    \begin{beamercolorbox}[wd=.34\paperwidth,ht=2.25ex,dp=1ex,center]{title in head/foot}%
      \usebeamerfont{title in head/foot}Experimental Analysis of Sorting Algorithms
    \end{beamercolorbox}%
    \begin{beamercolorbox}[wd=.33\paperwidth,ht=2.25ex,dp=1ex,right]{date in head/foot}%
      \usebeamerfont{date in head/foot}\insertframenumber{} / \inserttotalframenumber\hspace*{2ex}
    \end{beamercolorbox}}%
  \vskip0pt%
}

% --- Packages ---
\usepackage[T1]{fontenc}
\usepackage[utf8]{inputenc}
\usepackage{lmodern}
\usepackage{amsmath,amssymb}
\usepackage{booktabs}
\usepackage{pgfplots}
\pgfplotsset{compat=1.18}
\usepackage{array}
\usepackage{microtype}

% --- pgfplots cycle list: distinct colours + markers (B&W friendly) ---
\pgfplotscreateplotcyclelist{colormarkers}{%
  {blue!70!black,      mark=*,         thick},
  {red!70!black,       mark=square*,   thick},
  {green!50!black,     mark=triangle*, thick},
  {orange!80!black,    mark=diamond*,  thick},
  {violet,             mark=pentagon*, thick},
  {cyan!60!black,      mark=oplus*,    thick}%
}

% --- Metadata ---
\title[Sorting Algorithms]{Experimental Analysis of Sorting Algorithms}
\subtitle{An Empirical Study Across Sizes and Input Shapes}
\author{Maria Maiorescu}
\institute[UVT]{%
  Faculty of Informatics\\
  West University of Timi\c{s}oara
}
\date{\today}

% =====================================================================
\begin{document}

% =====================================================================
%  SLIDE 1 -- TITLE AND INTRODUCTION
% =====================================================================
\begin{frame}
  \titlepage
  \vspace{-0.5em}
  \begin{center}
    \small
    Goal: compare six classical sorting algorithms using different array sizes
    and input structures, combining theoretical complexity analysis with practical
    benchmarking implemented in~C.
  \end{center}
\end{frame}

% =====================================================================
%  SLIDE 2 -- IMPORTANCE OF SORTING
% =====================================================================
\begin{frame}{Importance of Sorting}
  \begin{columns}[T]
    \begin{column}{0.50\textwidth}
      \textbf{Sorting is fundamental in computer science:}
      \begin{itemize}
        \item Databases
        \item Search engines
        \item Operating systems
      \end{itemize}
      \vspace{0.7em}
      Theory focuses on asymptotic complexity, but practical performance also
      depends on \textbf{implementation details} and \textbf{input structure}.
    \end{column}
    \begin{column}{0.46\textwidth}
      \begin{alertblock}{Why benchmarking matters}
        Asymptotic bounds hide constants, cache effects, and the shape of real
        data.  An algorithm that wins on paper may lose on real input.
      \end{alertblock}
      \vspace{0.5em}
      \begin{block}{Consequence}
        Practical benchmarking is \emph{essential} alongside theoretical analysis.
      \end{block}
    \end{column}
  \end{columns}
\end{frame}

% =====================================================================
%  SLIDE 3 -- RELATED WORK
% =====================================================================
\begin{frame}{Related Work}
  \begin{itemize}
    \item \textbf{Donald E. Knuth (1998)} --- \textit{The Art of Computer Programming, Vol.~3}:
          foundational reference for sorting theory and complexity proofs.

    \vspace{0.4em}
    \item \textbf{Thomas H. Cormen, Charles E. Leiserson, Ronald L. Rivest,\\
          Clifford Stein (2022)} --- \textit{Introduction to Algorithms (CLRS)}:
          standard source for algorithm design and the $\Omega(n \log n)$ lower bound.

    \vspace{0.4em}
    \item \textbf{Nicolas Auger, Vincent Jug\'e, Cyril Nicaud, Carine Pivoteau (2019)}
          --- proved TimSort runs in $\mathcal{O}(n \log n)$ worst case; TimSort exploits
          adaptive behaviour of Insertion Sort on nearly-sorted data.

    \vspace{0.4em}
    \item \textbf{Robert Sedgewick \& Jon Bentley (2002)} --- argued that well-engineered
          Quick Sort achieves near-optimal average performance.
  \end{itemize}

  \vspace{0.5em}
  \begin{block}{What makes this project different}
    Unlike many previous studies, this project compares \textbf{six algorithms} under
    \textbf{identical testing conditions} and \textbf{multiple input structures}.
  \end{block}
\end{frame}

% =====================================================================
%  SLIDE 4 -- ALGORITHMS USED
% =====================================================================
\begin{frame}{Algorithms Used}
  \begin{columns}[T]
    \begin{column}{0.46\textwidth}
      \begin{itemize}
        \item \textbf{Bubble Sort} -- $\mathcal{O}(n^2)$, early-exit flag
        \item \textbf{Insertion Sort} -- $\mathcal{O}(n^2)$; drops to $\mathcal{O}(n)$
              on nearly-sorted input
        \item \textbf{Selection Sort} -- $\mathcal{O}(n^2)$, always
        \item \textbf{Merge Sort} -- $\mathcal{O}(n \log n)$, guaranteed
        \item \textbf{Quick Sort} -- $\mathcal{O}(n \log n)$ average;
              median-of-three pivot
        \item \textbf{Radix Sort (LSD)} -- $\mathcal{O}(d \cdot n)$,
              bypasses the comparison lower bound
      \end{itemize}
    \end{column}
    \begin{column}{0.52\textwidth}
      \centering
      \footnotesize
      \setlength{\tabcolsep}{3pt}
      \begin{tabular}{lcccc}
        \toprule
        \textbf{Algorithm} & \textbf{Best} & \textbf{Avg} & \textbf{Worst} & \textbf{Space} \\
        \midrule
        Bubble       & $\mathcal{O}(n)$       & $\mathcal{O}(n^2)$      & $\mathcal{O}(n^2)$      & $\mathcal{O}(1)$ \\
        Insertion    & $\mathcal{O}(n)$       & $\mathcal{O}(n^2)$      & $\mathcal{O}(n^2)$      & $\mathcal{O}(1)$ \\
        Selection    & $\mathcal{O}(n^2)$     & $\mathcal{O}(n^2)$      & $\mathcal{O}(n^2)$      & $\mathcal{O}(1)$ \\
        Merge        & $\mathcal{O}(n\lg n)$  & $\mathcal{O}(n \lg n)$  & $\mathcal{O}(n \lg n)$  & $\mathcal{O}(n)$ \\
        Quick        & $\mathcal{O}(n\lg n)$  & $\mathcal{O}(n \lg n)$  & $\mathcal{O}(n^2)$      & $\mathcal{O}(\lg n)$ \\
        Radix (LSD)  & $\mathcal{O}(dn)$      & $\mathcal{O}(dn)$       & $\mathcal{O}(dn)$       & $\mathcal{O}(n{+}d)$ \\
        \bottomrule
      \end{tabular}
      \vspace{0.5em}

      \textbf{Key takeaway:} the quadratic algorithms become inefficient on large
      arrays; Merge Sort and Quick Sort scale much better; Radix Sort achieved the
      fastest performance in most large tests.
    \end{column}
  \end{columns}
\end{frame}

% =====================================================================
%  SLIDE 5 -- BENCHMARK METHODOLOGY
% =====================================================================
\begin{frame}{Benchmark Methodology}
  \begin{columns}[T]
    \begin{column}{0.50\textwidth}
      \textbf{Benchmarking framework implemented entirely in C:}
      \begin{enumerate}
        \item \textbf{Data generators} --- fill array with the requested shape
        \item \textbf{Sorting routines} --- six in-place implementations
        \item \textbf{Timing layer} --- wrappers around \texttt{clock()}
        \item \textbf{Benchmarking loop} --- repeats runs, writes \texttt{results.csv}
      \end{enumerate}
      \vspace{0.5em}
      \textbf{Timing formula:}
      \[
        t = \frac{t_{\mathrm{end}} - t_{\mathrm{start}}}{\texttt{CLOCKS\_PER\_SEC}}
            \times 1000 \;\mathrm{ms}
      \]
    \end{column}
    \begin{column}{0.46\textwidth}
      \textbf{Input shapes tested:}
      \begin{itemize}
        \item \textbf{Random} arrays
        \item \textbf{Reversed} arrays
        \item \textbf{Partially sorted} at five levels:\\
              PS\,10\%, PS\,30\%, PS\,50\%, PS\,70\%, PS\,90\%
      \end{itemize}
      \vspace{0.5em}
      \begin{block}{Averaging strategy}
        Small-array tests ($n \leq 100$) are averaged over $10^6$ runs to
        reduce measurement noise.
      \end{block}
    \end{column}
  \end{columns}
\end{frame}

% =====================================================================
%  SLIDE 6 -- EXPERIMENTAL SETUP
% =====================================================================
\begin{frame}{Experimental Setup}
  \begin{columns}[T]
    \begin{column}{0.46\textwidth}
      \begin{block}{Compiler and hardware}
        Personal laptop, Unix-like OS\\
        \texttt{gcc -O2 -std=c99}
      \end{block}
      \vspace{0.6em}
      \textbf{Array sizes tested:}\\[0.3em]
      {\footnotesize$n \in \{10,\;50,\;100,\;1\,000,\;10\,000,\;100\,000,\;1\,000\,000\}$}
      Different input structures were used to observe how performance
      changes as the input size grows.
    \end{column}
    \begin{column}{0.50\textwidth}
      \textbf{Repetitions per measurement:}
      \begin{table}
        \centering
        \small
        \begin{tabular}{rl}
          \toprule
          \textbf{Size $n$}    & \textbf{Runs averaged} \\
          \midrule
          $10$, $50$, $100$    & $1{,}000{,}000$ \\
          $1{,}000$            & $100$ \\
          $10{,}000$           & $10$ \\
          $100{,}000$ and above & $1$ \\
          \bottomrule
        \end{tabular}
      \end{table}
      \vspace{0.4em}
      \textbf{Output:} \texttt{results.csv}\\
      Columns: \texttt{Elements, Runs, Algorithm, Variant, Elapsed~time(ms)}
    \end{column}
  \end{columns}
\end{frame}

% =====================================================================
%  SLIDE 7 -- RESULTS: SMALL AND MEDIUM ARRAYS
% =====================================================================
\begin{frame}{Results: Small and Medium Arrays}
  \begin{columns}[T]
    \begin{column}{0.48\textwidth}
      \centering
      \textbf{$n=100$ — avg.\ over $10^6$ runs (ms)}\\[0.3em]
      \tiny\setlength{\tabcolsep}{3pt}
      \begin{tabular}{lrrrrr}
        \toprule
        \textbf{Algorithm} & \textbf{Rev.} & \textbf{Rand.} & \textbf{PS10} & \textbf{PS50} & \textbf{PS90} \\
        \midrule
        Bubble    & 0.037 & 0.047 & 0.027 & 0.017 & 0.013 \\
        Quick     & 0.036 & 0.011 & 0.015 & 0.032 & 0.037 \\
        Insertion & 0.025 & 0.015 & 0.005 & 0.002 & 0.001 \\
        Selection & 0.023 & 0.031 & 0.026 & 0.026 & 0.026 \\
        Merge     & 0.010 & 0.016 & 0.012 & 0.010 & 0.010 \\
        Radix     & 0.005 & 0.007 & 0.005 & 0.005 & 0.005 \\
        \bottomrule
      \end{tabular}\\[0.6em]
      \textbf{$n=1{,}000$ — avg.\ over 100 runs (ms)}\\[0.3em]
      \begin{tabular}{lrrrrr}
        \toprule
        \textbf{Algorithm} & \textbf{Rev.} & \textbf{Rand.} & \textbf{PS10} & \textbf{PS50} & \textbf{PS90} \\
        \midrule
        Bubble    & 3.711 & 3.120 & 2.537 & 2.264 & 2.169 \\
        Quick     & 3.215 & 0.160 & 0.217 & 0.622 & 1.052 \\
        Insertion & 2.385 & 1.209 & 0.292 & 0.069 & 0.042 \\
        Selection & 2.093 & 2.388 & 2.347 & 2.345 & 2.363 \\
        Merge     & 0.124 & 0.214 & 0.150 & 0.135 & 0.130 \\
        Radix     & 0.066 & 0.103 & 0.087 & 0.065 & 0.067 \\
        \bottomrule
      \end{tabular}
    \end{column}
    \begin{column}{0.50\textwidth}
      \begin{tikzpicture}
        \begin{axis}[
          title={\footnotesize Random input},
          xlabel={\scriptsize Array size $n$},
          ylabel={\scriptsize Time (ms)},
          xmode=log, ymode=log,
          width=\linewidth, height=5.8cm,
          legend pos=north west,
          legend style={font=\tiny, cells={anchor=west}},
          grid=major,
          cycle list name=colormarkers,
          title style={font=\footnotesize},
          label style={font=\scriptsize},
          tick label style={font=\tiny}
        ]
          \addplot coordinates {(100,0.047)(1000,3.120)(10000,310.3)(100000,18521)};
          \addlegendentry{Bubble}
          \addplot coordinates {(100,0.011)(1000,0.160)(10000,2.1)(100000,10.4)};
          \addlegendentry{Quick}
          \addplot coordinates {(100,0.015)(1000,1.209)(10000,118.4)(100000,4199)};
          \addlegendentry{Insertion}
          \addplot coordinates {(100,0.031)(1000,2.388)(10000,230.3)(100000,6848)};
          \addlegendentry{Selection}
          \addplot coordinates {(100,0.016)(1000,0.214)(10000,2.8)(100000,28.2)};
          \addlegendentry{Merge}
          \addplot coordinates {(100,0.007)(1000,0.103)(10000,1.0)(100000,6.4)};
          \addlegendentry{Radix}
        \end{axis}
      \end{tikzpicture}
      \vspace{-0.3em}
      \scriptsize Radix Sort and Merge Sort already dominate at medium sizes.
      Insertion Sort becomes competitive when the input is almost sorted.
      Quadratic algorithms begin to slow down significantly.
    \end{column}
  \end{columns}
\end{frame}

% =====================================================================
%  SLIDE 8 -- RESULTS: LARGE ARRAYS
% =====================================================================
\begin{frame}{Results: Large Arrays}
  \begin{columns}[T]
    \begin{column}{0.40\textwidth}
      \centering
      \tiny\setlength{\tabcolsep}{2.5pt}
      \begin{tabular}{llrrr}
        \toprule
        $n$ & \textbf{Algo.} & \textbf{Rev.} & \textbf{Rand.} & \textbf{PS90} \\
        \midrule
        $10^4$  & Bubble    &  361.3 &  310.3 & 230.0 \\
                & Quick     &  312.1 &    2.1 &  10.6 \\
                & Insertion &  235.5 &  118.4 &   3.7 \\
                & Merge     &    1.5 &    2.8 &   1.7 \\
                & Radix     &    0.9 &    1.0 &   0.9 \\
        \midrule
        $10^5$  & Bubble    & 13\,228 & 18\,521 & 7\,640 \\
                & Quick     & 12\,726 &   10.4 &  72.1 \\
                & Insertion &  8\,343 & 4\,199 & 122.1 \\
                & Merge     &   20.6 &   28.2 &  21.4 \\
                & Radix     &    5.9 &    6.4 &   5.2 \\
        \midrule
        $10^6$  & Bubble    & \multicolumn{3}{c}{${\sim}10^6$\,ms (off chart)} \\
                & Quick     &  114.6 &  130.2 & 148.1 \\
                & Insertion & $812k$ & $397k$ & 11\,542 \\
                & Merge     &  232.7 &  359.7 & 227.0 \\
                & Radix     &   63.6 &   70.5 &  61.9 \\
        \bottomrule
      \end{tabular}
      \vspace{0.15em}
      
    \end{column}
    \begin{column}{0.58\textwidth}
      \begin{tikzpicture}
        \begin{axis}[
          title={\footnotesize Random input --- large $n$ (log-log)},
          xlabel={\scriptsize Array size $n$},
          ylabel={\scriptsize Time (ms)},
          xmode=log, ymode=log,
          width=\linewidth, height=5.8cm,
          legend pos=north west,
          legend style={font=\tiny, cells={anchor=west}},
          grid=major,
          cycle list name=colormarkers,
          title style={font=\footnotesize},
          label style={font=\scriptsize},
          tick label style={font=\tiny}
        ]
          \addplot coordinates {(10000,310.3)(100000,18521)(1000000,1808240)};
          \addlegendentry{Bubble Sort}
          \addplot coordinates {(10000,2.1)(100000,10.4)(1000000,130.2)};
          \addlegendentry{Quick Sort}
          \addplot coordinates {(10000,118.4)(100000,4199)(1000000,396592)};
          \addlegendentry{Insertion Sort}
          \addplot coordinates {(10000,230.3)(100000,6848)(1000000,657233)};
          \addlegendentry{Selection Sort}
          \addplot coordinates {(10000,2.8)(100000,28.2)(1000000,359.7)};
          \addlegendentry{Merge Sort}
          \addplot coordinates {(10000,1.0)(100000,6.4)(1000000,70.5)};
          \addlegendentry{Radix Sort}
        \end{axis}
      \end{tikzpicture}
    \end{column}
  \end{columns}
\end{frame}

% =====================================================================
%  SLIDE 9 -- ANALYSIS AND PROBLEMS ENCOUNTERED
% =====================================================================
\begin{frame}{Analysis and Problems Encountered}
  \begin{columns}[T]
    \begin{column}{0.50\textwidth}
      \begin{block}{Insertion Sort on nearly-sorted input}
        On PS\,90\% at $n=10^6$: \textbf{11\,542\,ms} vs.\ 396\,592\,ms on random
        data --- a $\mathbf{34\times}$ speedup.\\[0.3em]
        This is exactly the property that TimSort exploits.
      \end{block}
      \vspace{0.5em}
      \begin{block}{Quick Sort on structured inputs}
        Reversed $n=10{,}000$: \textbf{312.1\,ms} --- comparable to Bubble Sort.\\[0.3em]
        Even median-of-three produces unbalanced partitions on arithmetic sequences.
      \end{block}
    \end{column}
    \begin{column}{0.47\textwidth}
      \begin{alertblock}{Problem: Stack overflow}
        The recursive Quick Sort implementation caused a \textbf{segmentation fault}
        (stack overflow) on a partially-sorted $n=10^6$ array because of deep
        recursion.\\[0.3em]
        After adding the \textbf{median-of-three pivot} strategy the algorithm
        became much more stable.
      \end{alertblock}
      \vspace{0.5em}
      \begin{exampleblock}{Lesson}
        Practical implementations may behave very differently from theoretical
        expectations.
      \end{exampleblock}
    \end{column}
  \end{columns}
\end{frame}

% =====================================================================
%  SLIDE 10 -- CONCLUSION AND FUTURE WORK
% =====================================================================
\begin{frame}{Conclusion and Future Work}
  \begin{columns}[T]
    \begin{column}{0.50\textwidth}
      \textbf{Summary:}
      \begin{itemize}
        \item Compared \textbf{six classical sorting algorithms} across multiple
              array sizes and input structures.
        \item Confirmed the \textbf{limitations of quadratic algorithms}
              (unusable past $n \approx 10{,}000$).
        \item Highlighted the \textbf{efficiency of Merge Sort and Radix Sort}
              at scale.
      \end{itemize}
    \end{column}
    \begin{column}{0.47\textwidth}
      \textbf{Future work:}
      \begin{itemize}
        \item Test \textbf{additional data types} (floating-point, strings).
        \item Compare against \textbf{standard library sorts}
              (\texttt{qsort}, Python Timsort).
        \item Implement \textbf{iterative Quick Sort} to eliminate stack overflow.
        \item Profile \textbf{cache behaviour} with Cachegrind.
      \end{itemize}
    \end{column}
  \end{columns}

  \vspace{1em}
  \begin{center}
    \large\textbf{Thank you for your attention!}
  \end{center}
\end{frame}

\end{document}
